\chapter{Aplicaciones Industriales Avanzadas y de Ingeniería}
\label{cap:aplicaciones}

Los polímeros de ingeniería y funcionales avanzados han reemplazado a los metales estructurales en una gran variedad de industrias debido a su baja densidad, excelente resistencia a la corrosión y capacidad de sintonizar de manera precisa sus propiedades físicas y moleculares.

\section{Sustitución de Metales e Industria de Transporte}

\subsection{Aeroespacial}
En la aviación comercial moderna (por ejemplo, en el fuselaje y las alas del Boeing 787 Dreamliner y el Airbus A350 XWB), más del 50\% del peso estructural total de la aeronave está compuesto por materiales compuestos de matriz polimérica epóxica de alto rendimiento reforzada con fibra de carbono (CFRP). Esta reducción drástica de masa permite disminuir el consumo de combustible de las aerolíneas a nivel global.

\subsection{Automotriz}
El diseño de vehículos eléctricos modernos se apoya en el aligeramiento de peso (\textit{lightweighting}) para maximizar el rango de autonomía de las baterías de litio. Componentes mecánicos bajo el capó, colectores de admisión y soportes estructurales se moldean en poliamidas reforzadas con fibras de vidrio cortas.

\section{Polímeros en Biomedicina y Electrónica}

\subsection{Polímeros Conductores}
Históricamente, los plásticos se consideraban aislantes eléctricos perfectos. Sin embargo, en 1977 Shirakawa y colaboradores \cite{shirakawa1977} descubrieron que polímeros conjugados con enlaces dobles y simples alternados (como el poliacetileno dopado) exhiben una conductividad eléctrica comparable a la de los metales tradicionales. Esta propiedad sustenta hoy en día el diseño de células solares orgánicas (OPVs) y pantallas flexibles OLED comerciales.

\subsection{Barreras de Empaque}
En la conservación de alimentos y productos médicos, se diseñan películas multicapa coextruidas con polímeros especiales barrera (tales como el copolímero de etileno-vinil alcohol EVOH y el poli(cloruro de vinilideno) PVDC). Estos materiales presentan una permeabilidad al oxígeno extremadamente baja debido a su alto grado de polaridad y rigidez cristalina, retardando la oxidación bacteriana de los alimentos.
